% Suggested front matter matching the current claim-first manuscript framing.
% This file is not included by the build; papers/main.tex is authoritative.

% Suggested title
\title{Template Alignment Is Not Functional Alignment: Intervention-Budget Sensitivity in Small CNNs}

% Suggested abstract
\begin{abstract}
Interpretable-looking weights are often treated as evidence of interpretable
function, but weight-space resemblance and intervention behavior need not
agree. We test this gap in small CNNs with first-layer edge, corner, and ring
template priors and matched counterfactual interventions. Persistent
retention strongly preserves template alignment, yet it does not produce a
stable functional advantage. Instead, the release-versus-retention effect on
selected-channel fidelity changes sharply with intervention budget. In a
prospective confirmation on 20 fresh \texttt{two\_concepts} TinyCNN blocks,
the predeclared \emph{intervention-budget contrast}---whether release becomes
relatively more faithful at larger ($k=4,8$) than smaller ($k=1,2$) patch
budgets---is $B=0.33098$ (95\% CI $[0.27368,0.38827]$,
$p=2.28\times10^{-10}$). Release is strongly worse at $k=1,2$, slightly
worse at $k=4$, and slightly better at $k=8$. A separate 1,200-model
robustness study over 50 fresh blocks reproduces positive budget shifts for
TinyCNN on both tasks and across several prior schedules, whereas TwoLayerCNN
shows a different, often negative direction. The earlier positive
four-channel relative-usefulness effect is also partly baseline-driven,
showing that baseline choice can change the apparent functional conclusion.
An independent 32-checkpoint reimplementation audit passes all fixed
tolerances. Thus template alignment is not a sufficient summary of functional
influence in these models: the measured effect of a structured prior depends
on intervention scale, baseline definition, and architecture. We do not claim
a unique mechanism, human interpretability, or natural-image generalization.
\end{abstract}

% Suggested contribution paragraph
Our contributions are threefold. First, we make the gap between template
appearance and intervention behavior explicit under a controlled
structured-prior manipulation, while separately tracking prediction, concept
measurements, and learning dynamics. Second, we define and prospectively
confirm an intervention-budget contrast on fresh blocks, showing that the
release-versus-retention functional comparison changes materially with the
number of patched channels. Third, we map the boundary of that result across
1,200 additional models: the positive budget shift is robust in TinyCNN but
not architecture-invariant, and relative conclusions are also sensitive to
baseline construction. An independent checkpoint audit verifies the central
stored measurements within fixed numerical tolerances. The claims remain
restricted to the tested synthetic renderer, first-layer interventions, and
small CNN family.
